\scanpage{157}
$\mathcal Hom(M,N(\unsure{-j|j}))$ existe pour $j$ grand, dans la catégorie des motifs effectifs. Enfin, $\mathcal Hom(M,N)$ existe dans la catégorie de tous les motifs, comme on voit en prenant $\mathcal Hom(M,N(\unsure{j|-j}))(-j)$. On arrive ainsi à construire $\mathcal Hom(M,N)$ pour $M,N\in\operatorname{Ob}\widehat{\mathcal M}$ quelconques, avec les règles de calcul évidentes. On définit aussi $\check M=\mathcal Hom(M,\Lambda(0))$, et on trouve $\mathcal Hom(M,N)\simeq\check M\otimes N$. Mais on fera attention que si $M$ est effectif, $\check M$ ne l'est que si $\deg M=0$.

\[
B(M,N)\simeq\operatorname{Hom}(M\otimes N,\Lambda(-i))\simeq\operatorname{Hom}(M,N^{*})\simeq\operatorname{Hom}(N,M^{*})
\]
``formes bilinéaires'', avec $M,N$ de degré $i$.
Si $M=N$, on a la notion de forme alternée, forme symétrique.
\margin{\uline{N.B.} Forme non dégénérée en $M$ :}
\[
M\mathbin{-}N^{*}\quad\text{mono.}\quad\Longleftrightarrow\quad N\longrightarrow M^{*}\text{ est }\deleted{\text{injectif}}\inserted{\text{\unsure{épim.}}}
\]
On définit de même forme non dégénérée en $N$.

\margin{?}
Si $M=N$, les deux notions sont équivalentes car \deleted{tout hom} pour un hom $M\longrightarrow M^{*}$, dire que c'est \unsure{mono} ou \unsure{épimorphe} \deleted{\unsure{ou}} \unsure{un isom.} est pareil.
\mspage{156}
